\section{Pushdown automata}

\subsection{Why a stack?}
The language $A=\set{0^n1^n:n\geq 0}$ is not regular, so
no DFA recognizes it. The \emph{minimum extra memory} that does is a stack:
inspect the top, push, or pop.

Informal algorithm for $0^n1^n$:
\begin{itemize}
  \item while reading $0$s, push a marker;
  \item when $1$s begin, pop one marker per $1$;
  \item accept if the input ends with the stack empty (apart from a bottom
        marker).
\end{itemize}
A PDA has finitely many states but an unbounded stack. We allow
nondeterminism. Unlike finite automata, deterministic and nondeterministic
PDAs are \emph{not} equal in power. Nondeterministic PDAs are equivalent to
CFGs. Inherently ambiguous languages require nondeterminism; the converse is
false.

\subsection{Formal definition}
A (nondeterministic) PDA is a $6$-tuple.

\begin{definition}[PDA]
A \emph{pushdown automaton} is
\[
  P=(Q,\Sigma,\Gamma,\transition,q_0,F)
\]
where
\begin{enumerate}
  \item $Q$ is a finite set of states;
  \item $\Sigma$ is a finite input alphabet;
  \item $\Gamma$ is a finite stack alphabet;
  \item $\transition:Q\times\Sigma_\eps\times\Gamma_\eps\to\powerset{Q\times\Gamma_\eps}$
        is the transition function, with $\Sigma_\eps=\Sigma\cup\set{\eps}$ and
        $\Gamma_\eps=\Gamma\cup\set{\eps}$;
  \item $q_0\in Q$ is the start state;
  \item $F\subseteq Q$ is the set of accept states.
\end{enumerate}
\end{definition}

A transition $\bigl(q',b\bigr)\in\transition(q,a,x)$ means: in state $q$, read
$a$ from the input (or $\eps$: read nothing), pop $x$ from the stack (or
$x=\eps$: pop nothing), move to $q'$, and push $b$ (or $b=\eps$: push
nothing). Acceptance is by \emph{some} path: if any copy ends in an accept
state after consuming the input, the PDA accepts. A copy with no legal move
dies, as in an NFA.

\begin{definition}[Computation of a PDA]
Let $P=(Q,\Sigma,\Gamma,\transition,q_0,F)$ and $w\in\Sigma^*$. Write
$w=w_1\cdots w_m$ with $w_i\in\Sigma_\eps$. Then $P$ accepts $w$ if there are
states $r_0,\ldots,r_m\in Q$ and strings $s_0,s_1,\ldots,s_m\in\Gamma^*$
(stack contents, top on the left) such that
\begin{enumerate}
  \item $r_0=q_0$ and $s_0=\eps$ (or $s_0=\mathdollar$ if a bottom marker is pushed by
        an initial $\eps$-move);
  \item for each $i$, $\bigl(r_{i+1},b\bigr)\in\transition(r_i,w_{i+1},a)$
        where $s_i=at$ and $s_{i+1}=bt$ for some $a,b\in\Gamma_\eps$ and
        $t\in\Gamma^*$;
  \item $r_m\in F$.
\end{enumerate}
\end{definition}

\subsection{A PDA for \texorpdfstring{$0^n1^n$}{0n1n}}
\begin{example}
Let
\[
  Q=\set{q_{\mathrm{start}},q_{\mathrm{loop}},q_{\mathrm{pop}},q_{\mathrm{accept}}},\quad
  \Sigma=\set{0,1},\quad
  \Gamma=\set{0,\mathdollar},\quad
  F=\set{q_{\mathrm{accept}}}.
\]
\begin{itemize}
  \item From $q_{\mathrm{start}}$, $\eps$-move to $q_{\mathrm{loop}}$ pushing
        $\mathdollar$.
  \item In $q_{\mathrm{loop}}$ (have not yet seen a $1$): on $0$, push $0$.
  \item On $1$, pop a $0$ and go to $q_{\mathrm{pop}}$.
  \item In $q_{\mathrm{pop}}$: on $1$, pop a $0$; on $0$, no move (that copy
        dies).
  \item On $\eps$, if the top is $\mathdollar$, pop $\mathdollar$ and go to $q_{\mathrm{accept}}$.
\end{itemize}
\begin{center}
\begin{tikzpicture}[automaton,node distance=34mm]
  \node[state,initial] (qs) {$q_s$};
  \node[state,right=of qs] (ql) {$q_\ell$};
  \node[state,right=of ql] (qp) {$q_p$};
  \node[state,accepting,below=22mm of qp] (qa) {$q_a$};
  \path[->]
    (qs) edge node[align=center,above]{$\eps,\eps$\\$\to\mathdollar$} (ql)
    (ql) edge[loop above] node[align=center]{$0,\eps$\\$\to 0$} ()
         edge node[align=center,above]{$1,0$\\$\to\eps$} (qp)
    (qp) edge[loop above] node[align=center]{$1,0$\\$\to\eps$} ()
         edge node[align=center,right]{$\eps,\mathdollar$\\$\to\eps$} (qa);
\end{tikzpicture}
\end{center}
A label $a,x\to y$ means ``read $a$, pop $x$, push $y$''. On input $0011$,
after pushing $\mathdollar$ the stack is $00\mathdollar$ (top on the left);
each $1$ pops a $0$; a final $\eps$-move pops $\mathdollar$ and accepts.
\end{example}

\begin{example}[Balanced parentheses]
Let
\[
  L=\set{w\in\set{(,)}^*:w\text{ is balanced}}.
\]
Use states $q_s,q_\ell,q_a$, stack alphabet $\set{(,\mathdollar}$, and
\[
\begin{aligned}
  \transition(q_s,\eps,\eps)
    &=\set{(q_\ell,\mathdollar)},\\
  \transition(q_\ell,(,\eps)
    &=\set{(q_\ell,( )},\\
  \transition(q_\ell,),()
    &=\set{(q_\ell,\eps)},\\
  \transition(q_\ell,\eps,\mathdollar)
    &=\set{(q_a,\eps)}.
\end{aligned}
\]
Push on $($, pop on $)$, and accept by an $\eps$-move that pops the bottom
marker. Nondeterminism is not required. A copy dies if it tries to pop $($
when the top is $\mathdollar$, or if it reads $)$ with empty useful stack.
\end{example}

\subsection{PDAs and CFGs}
\begin{theorem}
A language is context-free if and only if some PDA recognizes it.
\end{theorem}

\begin{proof}[Proof idea]
($\Rightarrow$) From a CFG, a PDA keeps a sentential form on the stack, top
symbol being the next to expand. It nondeterministically replaces a variable
by a right-hand side, and matches terminals against the input.

($\Leftarrow$) From a PDA, a CFG has variables $A_{pq}$ meaning ``stack height
returns, going from state $p$ to state $q$''. Rules simulate a single PDA
move that pushes a symbol and later pops it, with an intermediate computation
in between.

Thus $\cfl=\set{\lang{P}:P\text{ is a PDA}}$.
\end{proof}

Acceptance by empty stack versus final state is equivalent for nondeterministic
PDAs: one extra state or marker converts either convention into the other.
